\chapter{Urobilinogen (Urine)}

\section*{What Is This Test?}
Urinary urobilinogen reflects the enterohepatic circulation of bilirubin. Conjugated bilirubin secreted into bile is converted by intestinal bacteria to colorless urobilinogen; most is re-excreted in feces (oxidized to urobilin giving feces their brown color), but a small fraction is reabsorbed, returned to the liver, and a portion filtered by kidneys into urine. Quantifying urinary urobilinogen helps distinguish jaundice causes: elevated in hepatocellular disease and hemolytic jaundice, absent in biliary obstruction. The test is part of the routine urinalysis dipstick.

\section*{Alternative Names}
\begin{itemize}
  \item Urine Urobilinogen
  \item Urinalysis Urobilinogen
  \item Urinary Urobilinogen
\end{itemize}
\section*{Principle}
Ehrlich's aldehyde reaction (modified) on the dipstick pad couples urobilinogen with p-dimethylaminobenzaldehyde to form a pink-red color. Sensitive at 0.2 mg/dL; semi-quantitative as normal, small (1+), moderate (2+), or large (3+). Quantitative 2-hour or 24-hour urine collection possible if needed.

\section*{History}
Ehrlich published the aldehyde reaction in 1901. Urobilinogen was first identified by Jaffe in 1869. Dipstick technology integrated urobilinogen detection in the 1950s.

\section*{Why This Test Is Ordered}
\begin{itemize}
  \item Differentiate types of jaundice: hepatocellular or hemolytic (elevated) vs biliary obstruction (absent)
  \item Workup of hemolytic anemia
  \item Hepatitis monitoring and chronic liver disease follow-up
  \item Standard urinalysis panel
\end{itemize}

\section*{Primary vs Secondary Test}
Primary as part of routine urinalysis; secondary in jaundice workup with serum bilirubin and liver function tests.

\section*{How This Test Is Performed}
A fresh random or 2-hour urine sample is collected. Reading requires fresh specimen; urobilinogen is photo-oxidized to urobilin with light exposure.

\section*{What Preparation Is Needed}
None. Antibiotics (especially broad-spectrum that suppress intestinal flora) reduce urobilinogen; ingestion of strict bowel rest may lower levels.

\section*{Contraindications and Precautions}
\begin{itemize}
  \item False positives: porphobilinogen (PBG) cross-reacts in acute intermittent porphyria; medications colored red (phenazopyridine), p-aminosalicylic acid, sulfonamides. False negatives: light exposure (labile urobilinogen), antibiotics (eliminate intestinal flora), severe cholestasis (little bilirubin reaches gut).
\end{itemize}
\section*{Risks}
None.

\section*{What Happens After the Test}
Results interpret in conjunction with urine bilirubin, serum bilirubin (total/direct), and liver enzymes.

\section*{Understanding Results}

\subsection*{Below Normal}
\begin{itemize}
  \item Biliary obstruction (choledocholithiasis, stricture, tumor): bilirubin reaches gut in insufficient quantities
  \item Broad-spectrum antibiotics suppressing intestinal flora
  \item Severe hepatocellular disease (decreased enterohepatic cycling)
\end{itemize}

\subsection*{Above Normal}
\begin{itemize}
  \item Hemolytic anemia: increased bilirubin production overwhelms hepatic conjugation; excess conjugated bilirubin reaches gut; urine urobilinogen rises
  \item Hepatocellular disease (hepatitis, cirrhosis, hepatotoxicity): impaired reuptake of reabsorbed urobilinogen
  \item Pulmonary or subdiaphragmatic hemorrhage resorption
  \item Constipation may transiently raise urobilinogen reabsorption
\end{itemize}

\section*{Normal Results}
\begin{itemize}
  \item \textbf{Urine urobilinogen (random dipstick):} 0.1--1.0 Ehrlich units/dL (normal or $<$ 1 mg/dL)
  \item \textbf{24-hour urine urobilinogen:} 0.5--4.0 mg/24 h
  \item \textbf{Clinical pattern interpretation:}
    \begin{itemize}
      \item High urine urobilinogen, normal urine bilirubin: hemolysis or hepatocellular disease
      \item Low urine urobilinogen, high urine bilirubin: biliary obstruction
      \item High urine urobilinogen, high urine bilirubin: hepatocellular disease
    \end{itemize}
\end{itemize}

\section*{Follow-up Tests}
\begin{itemize}
  \item \textbf{Serum bilirubin (total and direct)}: Fractionate hyperbilirubinemia
  \item \textbf{Liver function tests (ALT, AST, ALP, GGT)}: Distinguish hepatocellular vs cholestatic patterns
  \item \textbf{Urine bilirubin pad}: Classify jaundice (with urobilinogen)
  \item \textbf{Hemolysis panel (LDH, haptoglobin, reticulocyte count, peripheral smear)}: If hemolytic anemia suspected
  \item \textbf{Abdominal ultrasound, MRCP, or ERCP}: For biliary obstruction imaging
  \item \textbf{Urinalysis and urine culture}: For workup integration
\end{itemize}

\section*{References}
\begin{enumerate}
  \item Strasinger SK, Schumann GB. Urinalysis and Body Fluids. 7th ed. FA Davis; 2020.
  \item McPherson RA, Pincus MR, eds. Henry's Clinical Diagnosis and Management by Laboratory Methods. 24th ed. Elsevier; 2022.
  \item Wallach J. Interpretation Diagnostic Tests. 9th ed. Wolters Kluwer; 2022.
  \item Friedman LS, Keeffe EB. Handbook of Liver Disease. 4th ed. Elsevier; 2018.
\end{enumerate}

\clearpage
\section*{Regulatory Status and Authorization Date}
This chapter describes a test category, not a specific commercial product. FDA, CDSCO, EU IVDR, and MHRA approval or registration dates therefore cannot be assigned without the assay manufacturer, product name, instrument, and intended use. For laboratory-developed tests, an individual agency approval date may not exist. See the Preface for the interpretation of regulatory status.

\clearpage
